\section{A finite-bit signing interface}
\label{app:signing}

We record a conservative quantitative form of the rank-one signing
method of Ezeunala and Jiang~\cite{EzeunalaJiang2026}.  The constants
below suffice for the exact padding argument.  We include their
accounting and the finite-precision implementation to make the input
interface explicit.

\begin{theorem}[Rational rank-one signing]
\label{thm:rational-signing}
Given rational Hermitian matrices $H_1,\ldots,H_N$ of rank at most one,
a deterministic algorithm returns signs $\sigma_i\in\{-1,1\}$ such that
\[
 \left\|\sum_i\sigma_iH_i\right\|
 \le40\left\|\sum_iH_i^2\right\|^{1/2}.
\]
Its running time is polynomial in the total binary encoding length.
Consequently, for rational columns $v_i$ with $\sum_i v_iv_i^*=I$,
the same algorithm gives discrepancy at most
$40\sqrt{\max_i\|v_i\|^2}$.
\end{theorem}

The last assertion follows from
$(v_iv_i^*)^2=\|v_i\|^2v_iv_i^*$ and positivity.  Forming the outer
products has polynomial bit cost.  We prove the matrix assertion next.

\subsection{The potential and its quantitative estimates}

In dimension zero, any signs suffice.  Otherwise discard zero matrices;
if none remain, any signs again suffice.  A nonzero
rank-one Hermitian matrix has one nonzero eigenvalue, equal to its
trace.  With $D=\sum_i|\operatorname{Tr}H_i|$, put
\[
 A_i=D^{-1}\operatorname{sign}(\operatorname{Tr}H_i)H_i,
 \qquad \nu_0=\left\|\sum_iA_i^2\right\|.
\]
These are rational positive semidefinite matrices and
$\sum_i\operatorname{Tr}A_i=1$.  Matrix Cauchy--Schwarz gives
$1/(Nd^2)\le\nu_0\le1$, where $d$ is the matrix order.  Rational
positive-semidefiniteness tests and bisection compute rational parameters
\[
 \nu_0\le\nu\le\min\{1,1.001\nu_0\},\qquad
 \varepsilon=\nu/d,\qquad
 \frac{\sqrt\nu}{200N}\le\lambda\le\frac{\sqrt\nu}{100N}.
\]
Set $c=384$, $K=100$, and $\sigma=\lambda^2/\nu$.  Thus
$\varepsilon\ge1/(Nd^3)$ and
$1/(40000N^2)\le\sigma\le1/(10000N^2)$.

For $x\in[-1,1]^N$, let
\[
 S(x)=\sum_i x_iA_i,\quad
 \psi_i=(1-x_i^2)^{1/3},\quad
 \eta_x(Z)=c\sum_i\psi_iA_iZA_i.
\]
Define $R(x)$ as the minimum of
$t+\varepsilon\operatorname{Tr}(X+Y)$ over $X,Y\succ0$ subject to
\begin{equation}
 X^{-1}+S(x)+\eta_x(Y)\preceq tI,
 \qquad Y^{-1}-S(x)+\eta_x(X)\preceq tI,
 \label{eq:sign-potential}
\end{equation}
and put $\Psi(x)=R(x)+\lambda\sum_i\psi_i$.

The variational argument of Ezeunala and Jiang
\cite[Section~4.1]{EzeunalaJiang2026} uses only $c>0$ and gives
\begin{equation}
 \|S(x)\|\le R(x)\le\|S(x)\|+2\sqrt{(c+2)\nu},
 \qquad \Psi(0)\le(2\sqrt{386}+0.01)\sqrt\nu.
 \label{eq:sign-potential-bound}
\end{equation}
For clarity, its quantitative ingredients are as follows.  The minimizer
is unique; with its positive definite dual multipliers $P,Q$ it satisfies
\begin{gather}
 (2K)^{-1}I\preceq X,Y\preceq(K/\varepsilon)I,
 \quad P,Q\preceq I,\quad \operatorname{Tr}(P+Q)=1,\nonumber\\
 P=X\eta_x(Q)X+\varepsilon X^2,
 \qquad Q=Y\eta_x(P)Y+\varepsilon Y^2,
 \label{eq:sign-dual}
\end{gather}
and both inequalities in \eqref{eq:sign-potential} are tight.  Indeed,
the scalar choice $X=Y=aI$ proves the upper bound in
\eqref{eq:sign-potential-bound}; its value is below $42<K$.
On the sublevel set of objective at most $K$, the constraints bound the
inverse matrices and the trace penalty bounds $X,Y$.  Compactness gives
attainment.  Stationarity gives \eqref{eq:sign-dual}; positive dual
multipliers imply tightness.  Strict convexity of
$\operatorname{Tr}(PX^{-1})+\operatorname{Tr}(QY^{-1})$ gives uniqueness.

On pairs of Hermitian matrices define
\[
 \mathcal L(U,V)=
 (X^{-1}UX^{-1}-\eta_x(V),\,
  Y^{-1}VY^{-1}-\eta_x(U)).
\]
The positive-map argument of
\cite[Claim~4.1]{EzeunalaJiang2026} yields the following bound, with
the Frobenius norm specified explicitly:
\begin{equation}
 \|\mathcal L^{-1}\|\le\varepsilon^{-1},\qquad
 \mathcal L(P,Q)=\varepsilon(I,I).
 \label{eq:sign-L}
\end{equation}
All operator norms on spaces of matrix pairs here are induced by the
Frobenius norm. To verify the norm bound, write
$\mathcal L=\mathcal C(I-\mathcal T)$, where
\[
 \mathcal C(U,V)=(X^{-1}UX^{-1},Y^{-1}VY^{-1}),\qquad
 \mathcal T(U,V)=(X\eta_x(V)X,Y\eta_x(U)Y).
\]
Both $\mathcal T$ and $\mathcal C^{-1}$ preserve positive semidefinite
matrix pairs. The dual identities give
\[
 \mathcal T(P,Q)=(P,Q)-\varepsilon(X^2,Y^2).
\]
Since all four matrices are positive definite, there is a
$\delta>0$ such that the right side is at most $(1-\delta)(P,Q)$
in componentwise Loewner order. Thus $\mathcal T$ is a strict
contraction in the norm
\[
 \|(U,V)\|_{P,Q}
 =\inf\{s\ge0:-s(P,Q)\preceq(U,V)\preceq s(P,Q)\}.
\]
The convergent series
$\mathcal L^{-1}=\sum_{j\ge0}\mathcal T^j\mathcal C^{-1}$ is positive.
Its value on $(I,I)$ is $(P,Q)/\varepsilon\preceq(I,I)/\varepsilon$.
Consequently its induced norm for
$\|(U,V)\|_{\max}=\max\{\|U\|_{\op},\|V\|_{\op}\}$ is at most
$\varepsilon^{-1}$. Finally, $\mathcal L$ and its inverse are
self-adjoint for the Frobenius pairing. The Frobenius operator norm
of $\mathcal L^{-1}$ therefore equals its spectral radius, which is
bounded by every induced norm. This proves \eqref{eq:sign-L}.

The Hessian $\mathcal H$ of the Lagrangian in $(X,Y)$
satisfies
\begin{equation}
 \frac{2\varepsilon^2}{K}I\preceq\mathcal H,
 \qquad \|\mathcal H\|\le16K^3.
 \label{eq:sign-H}
\end{equation}
The lower bound follows by substituting $P\succeq\varepsilon X^2$ in
$2\operatorname{Tr}(PX^{-1}UX^{-1}UX^{-1})$ and using
$X^{-1}\succeq(\varepsilon/K)I$; the other block is identical.
The upper bound uses three inverse factors of norm at most $2K$ in
each of the two terms of a Hessian block.

\subsection{A quantified direction of decrease}

A coordinate is active when $|x_i|<1$; let $n$ be their number.
In the following differential calculations, vectors and derivatives
are restricted to these coordinates, with the other coordinates fixed
on the current open face.
Write $A_i=u_iu_i^*$ in this analysis and set
\[
 a_i=u_i^*Xu_i,\quad b_i=u_i^*Yu_i,\quad
 p_i=u_i^*Pu_i,\quad q_i=u_i^*Qu_i.
\]
The factors $u_i$ need not be rational: the algorithm below uses only
the rational matrices $A_i$.  If $c\psi_i b_i\ge1-x_i$, moving $x_i$
to $1$ preserves feasibility in \eqref{eq:sign-potential} with the
same $(t,X,Y)$ and decreases $\Psi$ by at least $\lambda\psi_i$.
The symmetric condition $c\psi_i a_i\ge1+x_i$ permits the move to $-1$.
We may therefore assume
\begin{equation}
 c\psi_i a_i<1+x_i,\qquad c\psi_i b_i<1-x_i
 \quad\text{for all active }i.
 \label{eq:sign-light}
\end{equation}

Put $\alpha_i=1+c\psi_i'b_i$, $\beta_i=1-c\psi_i'a_i$, and
\[
 T_i^2=\alpha_i\beta_i/\psi_i,\quad
 \ell_i=\alpha_i/T_i,\quad m_i=\beta_i/T_i,\qquad
 \widetilde a_i=\ell_i a_i,\quad \widetilde b_i=m_i b_i,\quad
 z_i=\ell_i p_i,\quad w_i=m_i q_i.
\]
For $k_i=-\psi_i''$, elementary substitution in
\eqref{eq:sign-light} gives
\begin{equation}
 \frac{k_i}{\alpha_i\beta_i}\ge\frac13,
 \quad\frac{k_i}{T_i^2}\ge\frac13,
 \quad T_i^{-2}\le\frac32,
 \quad\frac{|\psi_i'|}{\psi_i k_i}\le1,
 \quad\widetilde a_i\widetilde b_i\le c^{-2},
 \quad\widetilde a_i^2,\widetilde b_i^2\le3c^{-2}.
 \label{eq:sign-scalar}
\end{equation}
Here is a direct check.  By symmetry take $x_i\ge0$ and put
$r=1-x_i^2$.  Then
$2/3\le\alpha_i\beta_i\le2/r$ and
$k_i=(2/3)r^{-5/3}(1+x_i^2/3)$.
Also $\widetilde a_i\widetilde b_i=\psi_i a_i b_i
\le r^{2/3}/c^2$.  For
$u=c\psi_i a_i/(1+x_i)$ and $v=c\psi_i b_i/(1-x_i)$ in $(0,1)$,
\[
 c^2\widetilde a_i^2
 =r^{2/3}u^2\frac{1+x_i-2x_iv/3}{1-x_i+2x_iu/3},\qquad
 c^2\widetilde b_i^2
 =r^{2/3}v^2\frac{1-x_i+2x_iu/3}{1+x_i-2x_iv/3}.
\]
Maximizing the first in $u$ and minimizing $v$, and maximizing the
second in $u,v$, bounds them by $3$ and $1$.
For negative $x_i$ the roles interchange.  These formulas also prove
the remaining assertions of \eqref{eq:sign-scalar}.

We quantify the covariance construction in
\cite[Lemma~4.5]{EzeunalaJiang2026}.  Define the real symmetric,
entrywise nonnegative positive semidefinite matrices
\[
 \mathsf A_{ij}=c\ell_i\ell_j|u_i^*Xu_j|^2,
 \qquad \mathsf B_{ij}=cm_im_j|u_i^*Yu_j|^2,
\]
and $D_F=\operatorname{diag}(\widetilde a_i z_i)$,
$D_G=\operatorname{diag}(\widetilde b_i w_i)$.
The dual identities give $\mathsf A w\le z$ and
$\mathsf Bz\le w$ coordinatewise.  Moreover,
\[
 \|D_G^{1/2}\mathsf A D_F^{-1/2}\|_F^2\le n/c,
 \qquad
 \|D_F^{1/2}\mathsf B D_G^{-1/2}\|_F^2\le n/c.
\]
For example, $\mathsf A_{ij}^2\le
c\widetilde a_i\widetilde a_j\mathsf A_{ij}$ and
$\widetilde a_i\widetilde b_i\le c^{-2}$ bound the first sum by
$c^{-1}\sum_j(\mathsf A w)_j/z_j$.

Restrict $D_F^{1/2}F$ and $D_G^{1/2}G$ to right singular directions
with singular value at most $1/2$ for these two maps.  Each restriction
has codimension at most $4n/c$.  Impose also
$(I-\mathsf A)F+(I-\mathsf B)G=0$.  The resulting subspace
$\mathcal G\subseteq\mathbb R^{2n}$ has dimension at least $n-8n/c$.
The covariance theorem of Bansal and Garg~\cite[Theorem~6]{BansalGarg2017},
in the form stated by Bansal~\cite[Theorem~2.5]{Bansal2024} and recalled
in \cite[Lemma~2.1]{EzeunalaJiang2026}, supplies a
positive semidefinite $\Gamma$ with
\[
 \operatorname{range}\Gamma\subseteq\mathcal G,\quad
 \Gamma\preceq4\operatorname{diag}\Gamma,\quad
 \Gamma_{jj}\le1,\quad
 \operatorname{Tr}\Gamma\ge(1/2-8/c)n=23n/48.
\]
Take a centered finitely supported $(F,G)$ with covariance $\Gamma$
and set $H=F-\mathsf BG=\mathsf AF-G$, $h_i=H_i/T_i$, and
$\omega_i=\mathbb E H_i^2$.
The two singular-value restrictions give
\[
 \|F\|_{D_F}^2+\|G\|_{D_G}^2
 \le4\sum_i(\widetilde a_i z_i+\widetilde b_i w_i)H_i^2.
\]
The response matrices are
$\dot X=X(\sum_i\ell_iF_iA_i)X$ and
$\dot Y=Y(\sum_i m_iG_iA_i)Y$.
The diagonal of their inversion-energy quadratic form
\[
 E_h=\operatorname{Tr}(PX^{-1}\dot X X^{-1}\dot X X^{-1})
     +\operatorname{Tr}(QY^{-1}\dot Y Y^{-1}\dot Y Y^{-1})
\]
is $(D_F,D_G)$.  Hence
\begin{equation}
 \mathbb E E_h\le16\sum_i
          (\widetilde a_i z_i+\widetilde b_i w_i)\omega_i.
 \label{eq:sign-energy}
\end{equation}

We expand the first and second variation calculation of
\cite[Lemma~4.3 and its proof]{EzeunalaJiang2026} to obtain
\[
 \partial_iR=T_i(z_i-w_i),\qquad
 h^T\nabla^2Rh\le3E_h-\frac c2\sum_i k_i h_i^2(b_ip_i+a_iq_i).
\]
The first identity is the envelope formula
$\partial_iR=\alpha_i p_i-\beta_i q_i$.
For the second, keep $t$ fixed and replace $x$ by $x+\tau h$.
Invertibility of $\mathcal L$ and the implicit function theorem
give a local curve $(X(\tau),Y(\tau))$ satisfying both tight
constraints. Its first derivatives are the response matrices above:
the identities $H=F-\mathsf BG=\mathsf AF-G$ verify the linearized
constraints. This curve stays positive definite for small $\tau$,
and its objective bounds $R(x+\tau h)$ from above with equality at
$\tau=0$. Differentiating twice and pairing with $P,Q$, using
$\mathcal L(P,Q)=\varepsilon(I,I)$, therefore gives
\begin{align*}
 h^T\nabla^2Rh
 &\le2E_h-c\sum_i k_i h_i^2(b_ip_i+a_iq_i)\\
 &\quad+2c\sum_i\psi_i'h_i
       \bigl(q_i u_i^*\dot X u_i+p_i u_i^*\dot Y u_i\bigr).
\end{align*}
Write $E_X,E_Y$ for the two summands of $E_h$. Since
$P\succeq X\eta_x(Q)X$,
\[
 E_X\ge\operatorname{Tr}(\eta_x(Q)\dot X X^{-1}\dot X)
     =c\sum_i\psi_iq_i u_i^*\dot X X^{-1}\dot X u_i.
\]
The inequality
$|u_i^*\dot X u_i|^2\le a_i u_i^*\dot X X^{-1}\dot X u_i$
and weighted Cauchy--Schwarz bound the $X$ mixed sum by
\[
 \left|c\sum_i\psi_i'h_iq_i u_i^*\dot X u_i\right|
 \le\sqrt{E_X B_X},\qquad
 B_X=c\sum_i(\psi_i')^2\psi_i^{-1}h_i^2a_iq_i.
\]
The corresponding $Y$ bound uses
$B_Y=c\sum_i(\psi_i')^2\psi_i^{-1}h_i^2b_ip_i$.
Thus the full mixed term is at most $2\sqrt{E_hB_h}$, where
$B_h=B_X+B_Y$. Finally,
$(\psi_i')^2/(\psi_i k_i)\le1/2$ gives
$B_h\le(c/2)\sum_i k_i h_i^2(b_ip_i+a_iq_i)$, and
$2\sqrt{E_hB_h}\le E_h+B_h$ proves the asserted Hessian bound.

Let $\Sigma=\mathbb E hh^T$ and choose
$\mu_i=-24T_i^{-1}(\widetilde a_i-\widetilde b_i)\omega_i$.
Using \eqref{eq:sign-scalar} and \eqref{eq:sign-energy}, the operator
$\mathcal Df=\nabla f\cdot\mu+\tfrac12\operatorname{Tr}(\nabla^2f\Sigma)$
satisfies
\begin{align*}
 \mathcal DR
 &\le(24-c/12)\sum_i
       (\widetilde a_i w_i+\widetilde b_i z_i)\omega_i\le0,\\
 \mathcal D\Bigl(\sum_i\psi_i\Bigr)
 &\le-\left(\frac12-\frac{48}{c}\right)
                  \sum_i\frac{k_i\omega_i}{T_i^2}
 \le-\frac18\sum_i\omega_i.
\end{align*}
The first line uses cancellation of
$24\sum_i(\widetilde a_i z_i+\widetilde b_i w_i)\omega_i$ by the drift.
For the second, the identity $T_i(\widetilde a_i+\widetilde b_i)=a_i+b_i$
and \eqref{eq:sign-light} give
$|\psi_i'\mu_i|\le(48/c)k_i\omega_i/T_i^2$.

All movement bounds are polynomial and uniform in the state.
Indeed \eqref{eq:sign-dual} gives
\[
 \mathsf Aw\le(1-s_0)z,\quad
 \mathsf Bz\le(1-s_0)w,\quad
 s_0=\varepsilon/(4K^2),\quad
 \|\mathsf A\|,\|\mathsf B\|\le3N.
\]
The spectral radius of $\mathsf A\mathsf B$ is at most $(1-s_0)^2$.
Using its positive semidefinite conjugate yields
\[
 \|(I-\mathsf A\mathsf B)^{-1}\|,
 \|(I-\mathsf B\mathsf A)^{-1}\|
 \le J:=1+36K^2N^2/\varepsilon.
\]
Solving the two response equations gives
$\|F\|^2+\|G\|^2\le E_0\|H\|^2$, where
$E_0=2J^2(1+3N)^2$.  Consequently, with
\[
 C_1=\max\{1,48E_0/23,4(1+9N^2)\},\qquad C_0=8C_1,
\]
we have $\sum_i\omega_i\ge n/C_1$, $\omega_i\le C_1$ and
\begin{equation}
 \|\mu\|_\infty\le C_0,\quad
 \operatorname{Tr}\Sigma\le C_0n,\quad
 \mathcal D\Psi\le-\lambda n/C_0.
 \label{eq:sign-progress}
\end{equation}
This proves a quantitative descent alternative for $c=384$.
The covariance and drift are used only to prove that a direction of
decrease exists.  The algorithm below computes neither of them nor
the possibly irrational spectral subspaces used in their construction.

Here is the deterministic alternative that the algorithm will use.
Suppose at least one coordinate is active, every active coordinate is
farther than $\sigma$ from its endpoints, and put
$a=\lambda/(8C_0^2)$. Either an endpoint test above
succeeds, giving a decrease of at least $\lambda\sigma$, or
\begin{equation}
 \|\nabla\Psi(x)\|_\infty\ge a
 \quad\text{or}\quad
 \lambda_{\min}(\nabla^2\Psi(x))\le-8a.
 \label{eq:sign-deterministic-alternative}
\end{equation}
The endpoint estimate uses $\psi_i\ge\sigma$. If all endpoint tests
fail, \eqref{eq:sign-light} holds and the preceding calculation
applies. If also $\|\nabla\Psi\|_\infty<a$, then
$\nabla\Psi\cdot\mu\ge-aC_0n$, whereas
$\mathcal D\Psi\le-\lambda n/C_0=-8aC_0n$. Hence
\[
 \tfrac12\operatorname{Tr}(\nabla^2\Psi\,\Sigma)
       \le-7aC_0n.
\]
Together with $\operatorname{Tr}\Sigma\le C_0n$, this gives a
Hessian eigenvalue at most $-14a$, and therefore the asserted
$-8a$ bound. The remaining task is to find one of these decreases
using rational value queries.

\subsection{Uniform derivatives and a rational value oracle}

We next make the descent computable.  On any open face, let
$z=(t,X,Y,P,Q)$ and write its tight constraints and stationarity
equations as $F(x,z)=0$.  The source's smoothness proof
\cite[Section~4.1]{EzeunalaJiang2026} applies to our parameters.
Here we record a uniform quantitative bound.
Write $B=\mathcal L^{-1}$ and $E=(I,I)$.  An inhomogeneous
linearized system has the form
\[
 \mathcal Lu+hE=f,\qquad
 \mathcal Lp-\mathcal Hu=g,\qquad \langle E,p\rangle=r.
\]
Eliminating $u,p$ determines $h$ with denominator
$\langle BE,\mathcal HBE\rangle$.
By \eqref{eq:sign-L}, $BE=(P,Q)/\varepsilon$; thus
\[
 \|BE\|_F\le\varepsilon^{-1},\qquad
 \langle BE,\mathcal HBE\rangle
 \ge\frac{2\varepsilon^2}{K}\|BE\|_F^2
 \ge\frac1{Kd}.
\]
The last inequality uses
$\|P\|_F^2+\|Q\|_F^2\ge1/(2d)$.
The upper bound uses $\operatorname{Tr}(P+Q)=1$ and positivity.
Back substitution in the three equations and
\eqref{eq:sign-L}--\eqref{eq:sign-H} bound the complete Jacobian
inverse by a constant times $(d+1)(1+\varepsilon^{-1})^4$.

When every active coordinate is at least $\sigma/2$ from its endpoints,
the first three derivatives of $\psi_i$ are bounded by
$100(2/\sigma)^3$.  All derivatives of $F$ through order three are
therefore bounded by a computable polynomial in
$N,d,\varepsilon^{-1},\sigma^{-1}$.
Three implicit differentiations give a uniform polynomial bound $M\ge1$
for both $\|\nabla^2\Psi\|$ and $\|\nabla^3\Psi\|$ on these strips.
There is no dependence on the number of faces.  For an explicit
choice, let
\begin{gather*}
 A=10^{18}(d+1)(1+\varepsilon^{-1})^4,\\
 D_*=10^{20}(N+1)^3(d+1)^3
             (1+\varepsilon^{-1})^3(1+\sigma^{-1})^3,\\
 Z_1=AD_*,\quad Z_2=AD_*(1+Z_1)^2,\quad
 Z_3=AD_*\bigl((1+Z_1)^3+3(1+Z_1)Z_2\bigr),\\
 M=\left\lceil1+(1+2d\varepsilon)(Z_2+Z_3)
                         +800\lambda N\sigma^{-3}\right\rceil.
\end{gather*}
The product rule, the inverse bound $2K$, and Frobenius product
inequalities verify these deliberately loose bounds.

A value of $\Psi$ to additive accuracy $\eta>0$ is computable in
polynomial time in the input length and $\log\eta^{-1}$.
Approximate every $\psi_i$ by a nonnegative rational in $[0,1]$ to
error $\delta$.  This changes the optimal value by at most
$(cKd+\lambda N)\delta$: in either direction, increase $t$ by
$cKd\delta$ to absorb the perturbation of the matrix constraints.
Indeed, $\sum_iA_i^2\preceq\nu I$ and
$\|Y\|\le K/\varepsilon$ bound the change in $\eta_x(Y)$ by
$c\delta\nu\|Y\|\le c\delta\nu K/\varepsilon=cKd\delta$;
the other block is identical.
The resulting rational program is a semidefinite program by the
Schur-complement formula.  Impose the rational bounds
\[
 I/400\preceq X,Y\preceq(200/\varepsilon)I,
 \qquad 0\le t\le100.
\]
The scalar upper bound and sublevel estimates hold for every coefficient
vector in $[0,1]^N$, so these boxes contain the minimizers both before
and after approximation. The preceding comparison of unboxed optimal
values therefore applies to the boxed programs as well.
The point $X=Y=I/20$, $t=50$ has
constraint slack at least $50-20-2-384/20=8.8$ and lies strictly
inside the bounds.  It supplies an inverse-polynomial interior
radius: each Schur block dominates
$\left(\begin{smallmatrix}28.8I&I\\I&0.05I\end{smallmatrix}\right)
\succeq I/100$, since subtracting $I/100$ leaves the scalar
determinant $28.79\cdot0.04-1=0.1516>0$.
In upper-triangular real and imaginary coordinates, a ball of radius
$10^{-6}$ preserves this margin and all the boxes: a matrix-pair
perturbation changes either Schur block by at most
$|\Delta t|+384\|\Delta Y\|_F+\|\Delta X\|_F$ or its symmetric
counterpart.  An outer radius $400(d+1)/\varepsilon$ suffices.
Rational weak optimization~\cite{GLS1988} therefore evaluates the
rational program to error $\eta/2$ in polynomial work.  Choosing
$\delta\le\eta/[2(cKd+\lambda N)]$ gives the claimed oracle.
Rational bisection computes the cube roots.  This value oracle also
works on lower-dimensional faces and at vertices.

\subsection{Rational descent and termination}

At the start of an iteration, move every active coordinate at distance
$z\le\sigma$ from its nearest endpoint to that endpoint.  This never
increases the exact potential, since
\[
 \Delta\Psi\le z\sqrt\nu-\lambda[z(2-z)]^{1/3}
             \le z\sqrt\nu-\lambda\sqrt z\le0.
\]
There are at most $N$ such permanent freezes.  Every remaining active
coordinate is more than $\sigma$ from either endpoint. If none remain,
return the resulting vertex; all derivative and direction computations
below are performed only when $n\ge1$.

With $a=\lambda/(8C_0^2)$ as above, choose positive powers of two $s,h$, each within
a factor two below its indicated bound, with
\[
 s\le\min\{\sigma/4,a/(4M),1/4\},\qquad
 h\le\min\{\sigma/8,1,a/(64MN)\},\qquad
 \Delta=as^2.
\]
Central first differences, central diagonal second differences, and
four-corner mixed second differences of the value oracle, each query
to error at most $ah^2/(128N)$, give rational approximations
$\widehat g,\widehat H$ to $\nabla\Psi,\nabla^2\Psi$ satisfying
\[
 \|\widehat g-\nabla\Psi\|_\infty<a/8,
 \qquad\|\widehat H-\nabla^2\Psi\|<a/16.
\]
Taylor's theorem and the third-derivative bound prove these estimates;
all stencil points stay in the $\sigma/2$ strips.

Consider all $2n$ single-coordinate endpoint moves and the coordinate
move $x-s\operatorname{sign}(\widehat g_i)e_i$, where
$|\widehat g_i|$ is maximal.  Test the rational matrix
$\widehat H+7aI$ for positive semidefiniteness by exact elimination.
If it fails, elimination supplies a nonzero rational vector of
negative quadratic form.  Set $\beta=a/(4096M)$ and
$\tau=a/(1024M)$ and normalize that vector, using a rational upper
approximation to its norm, to obtain $w$ with
$1-\beta\le\|w\|\le1$.  Its Rayleigh quotient for
$\widehat H$ is strictly below $-7a$.
Choose $\rho$ as the largest dyadic power at most
$sa/(4096MN)$; it divides the dyadic $s$.
Round $s(1-\tau)w$ to $\rho\mathbb Z^n$ to obtain $q$.
The rounding error is at most $\sqrt n\rho/2\le s\tau/8$, so
$\|q\|\le s$ and $\|q-sw\|\le9s\tau/8$.
The negative-direction certificate and the Hessian error bound give
\[
 w^T\nabla^2\Psi\,w<-\frac{111}{16}a(1-\beta)^2,
 \qquad
 |q^T\nabla^2\Psi\,q-s^2w^T\nabla^2\Psi\,w|
 \le\frac9{4096}as^2.
\]
In particular,
\[
 q^T\nabla^2\Psi\,q\le-6as^2.
\]
Include both $x+q$ and $x-q$.  If the positive-semidefiniteness test
passes, omit these two candidates.  Whenever
$\lambda_{\min}(\nabla^2\Psi)\le-8a$, the test necessarily fails,
so a curvature candidate is available when it is needed.
The two displacements are exact negatives, so
their linear Taylor terms cancel.

By \eqref{eq:sign-deterministic-alternative}, one candidate decreases
the true potential by at least $2\Delta$.
An available endpoint move gives a decrease at least
$\lambda\sigma\ge2\Delta$. In the remaining case, if
$\|\nabla\Psi\|_\infty\ge a$, the selected coordinate has the
correct sign and derivative magnitude at least $3a/4$, giving decrease
at least $(3/4)as-Ms^2/2\ge(5/8)as\ge(5/2)\Delta$.
Otherwise \eqref{eq:sign-deterministic-alternative} gives
$\lambda_{\min}(\nabla^2\Psi)\le-8a$; averaging the two curvature
candidates gives decrease at least
$3as^2-Ms^3/6>2\Delta$.
Evaluate all final candidates to error $\Delta/4$ and choose the
smallest reported value.  Selection loses at most $\Delta/2$ against
the best true candidate.  Every accepted step thus decreases the
exact potential by at least $3\Delta/2$.

Starting from $x=0$, every live coordinate stays on the fixed mesh
$\rho\mathbb Z$; frozen coordinates equal $\pm1$.  All coordinates,
parameters, and query points have polynomial encoding length.
There are at most $\Psi(0)/\Delta$ main iterations, each with
polynomially many rational semidefinite value queries and one exact
rational positive-semidefiniteness test. The normalized lower bounds
on $\nu$, $\varepsilon$, $\lambda$, and $\sigma$ imply that $C_0$ and
$M$ are bounded by fixed polynomials in $N+d$. In turn, $a$ and the
chosen $s,h$ are bounded below by fixed inverse polynomials, as is
$\Delta=as^2$. Since $\Psi(0)<40$, the iteration count is polynomial
as well. Together with the encoding and oracle bounds, this proves
deterministic polynomial bit complexity on every execution.

At the final vertex $x^f$, the coordinate term vanishes and
\eqref{eq:sign-potential-bound} gives
\[
 \left\|\sum_i x_i^f A_i\right\|
 \le(2\sqrt{386}+0.01)\sqrt\nu
 \le(2\sqrt{386}+0.01)\sqrt{1.001\nu_0}
 <40\sqrt{\nu_0}.
\]
Undoing the rational normalization and folding the signs of the
original traces into $x_i^f$ proves
Theorem~\ref{thm:rational-signing}.
